\track{T1\quad 对偶、群论与耦合常数}{本轨道 14 条：先把「群怎么对偶、耦合怎么变换、通量取值于谁」钉死，这是全书地基。}

% ============================================================
% 1. 电-磁对偶
% ============================================================
\begin{word}{电-磁对偶}{electric--magnetic duality}{elec-mag-dual}{§1 p.2}
\mathdef{$4$ 维流形上 $2$-形式方程组的对称性：闭形式方程 $dF=0$ 与余闭方程 $d\star F=0$ 在 $F\mapsto\star F$ 下互换（$\star$ 为 Hodge 星，$\star^2=\pm1$，故这是方程本身的精确对称）。两个不变量分别是对称的拓扑量 $\int F\wedge F$ 与能量 $\int F\wedge\star F$。非阿贝尔化之后被互换的不再是场而是电荷格：电荷取值于权格 $\Lambda_{\mathrm{wt}}$，磁荷取值于其对偶格 $\Lambda_{\mathrm{cwt}}$，两格之间有完美配对 $\langle n,m\rangle\in\mathbb Z$。}
\mathex{取 $G=U(1)$、底流形为 $\mathbb R^4$。$F\mapsto\star F$ 把 $(E,B)$ 送到 $(B,-E)$，保持 $E^2-B^2$ 与 $E\cdot B$ 不变。在一根无穷长直线电荷外取绕它的 $S^2$：$\int_{S^2}\star F=4\pi q$；对偶后在同一个 $S^2$ 上 $\int_{S^2}F=4\pi g$，且 $qg\in2\pi\mathbb Z$。}
\physdef{$4$ 维 $\mathcal N=4$ 超对称杨--米尔斯理论中的精确量子对称：把规范粒子与磁单极子互换、耦合 $e\mapsto4\pi/e$、规范群 $G\mapsto{}^LG$。它同时改三样东西：场强、电荷格、群的整体形式。在 $\mathcal N=4$ 里它不只是一个弱--强耦合的近似陈述。}
\physex{$G=SU(2)$、$\theta=0$：$W^\pm$ 玻色子质量约为 $ev$，磁单极子质量约为 $(4\pi/e)v$。$e\ll1$ 时后者极重，这正是弱耦合下看不见磁单极子的原因；$S$-对偶交换两套谱，于是强耦合的 $SU(2)$ 理论可以用弱耦合的 $SO(3)$ 理论来算。}
\bridge{§1 p.2--3：物理侧的电-磁对偶群就是根数据对偶化所得的 ${}^LG$，耦合按 $\tau\mapsto-1/n_{\mathfrak g}\tau$ 变换（§2.2）。这是全书地基：先把换群规则钉死，§3 的扭化与 §4--§5 的镜像对称才有明确的 $(G,{}^LG)$ 对可言。}
\intuit{直觉上像傅里叶变换：电变量（$A_\mu$、Wilson 线）与磁变量（对偶规范场、't Hooft 线）互换而物理不变。这个类比靠得住，是因为 $\mathcal N=4$ 的超共形对称排除了大部分量子修正，使经典对偶成为精确对偶。}
\pitfall{按字面把电-磁对偶理解成「$E$ 与 $B$ 互换」会推出「非阿贝尔情形也一样」；错在漏掉换群这一步：$G$ 与 ${}^LG$ 一般不同构（$SU(N)$ 的对偶是 $PSU(N)$），只换场不换群描述的是另一个理论。}
\end{word}

% ============================================================
% 2. Montonen--Olive 对偶
% ============================================================
\begin{word}{Montonen--Olive 对偶}{Montonen--Olive duality}{montonen-olive}{§1 p.2}
\mathdef{固定一个单李代数 $\mathfrak g$，把「电荷格与磁荷格互换、耦合取倒数」的操作组织成一个陈述：以 $\mathfrak g$ 为 Lie 代数的规范理论在参数 $\tau$ 处的谱（规范粒子与磁单极子的质量、电荷）与它在 $-1/\tau$ 处的谱相同，只是电与磁的角色互换。其群论内容是根数据的对偶化：$G\mapsto{}^LG$，权格与余权格互换。原始猜想只含一个 $\mathbb Z_2$（无 $\theta$ 角）。}
\mathex{$G=SU(2)$、$\theta=0$：$\tau=4\pi i/e^2$，猜想断言理论与自身对偶，耦合 $e\mapsto4\pi/e$。自对偶点由 $e^2=4\pi$ 给出，此时 $\tau=i$、$\Im\tau=1$。加上 $\theta$ 后 $\mathbb Z_2$ 扩为 $SL(2,\mathbb Z)$：$\tau\mapsto\tau+1$ 与 $\tau\mapsto-1/\tau$ 都成为对称。}
\physdef{$4$ 维 $\mathcal N=4$ 超对称杨--米尔斯理论中「存在一个等价描述，其规范群为 ${}^LG$、耦合为 $4\pi/e$」的猜想。物理内容是谱的互换：基本粒子 $\leftrightarrow$ 't Hooft--Polyakov 单极子，$W$ 玻色子 $\leftrightarrow$ 磁单极子，Wilson 线 $\leftrightarrow$ 't Hooft 线。它只在超对称理论里才可能成立。}
\physex{$G=SU(2)$、$e^2=4\pi$（即 $\tau=i$）：理论与自身的对偶描述重合，$W$ 玻色子与磁单极子质量都等于 $\sqrt{4\pi}\,v$。$e$ 偏离这个值时两侧耦合不同，但 BPS 谱仍逐项对上：第 $n$ 个电荷扇区的质量公式在 $e\mapsto4\pi/e$ 下不变。}
\bridge{§1 p.2：KW06 把 Montonen--Olive 猜想当作出发点，真正的工具是它的精确版本 $S$-对偶（§2.2）。注意原猜想只含 $g\leftrightarrow1/g$ 的 $\mathbb Z_2$；$\theta$ 的加入使它与非单连通型的 lacing 数一起生成 $SL(2,\mathbb Z)$ 或其子群，这才是 §2.2 的内容。}
\intuit{直觉是「重与轻互换」：强耦合下重的对象在弱耦合描述里变轻，于是难算的理论变成好算的理论。靠得住的原因不是微扰论，而是超对称使 BPS 态的质量由中心荷精确给出，从而可以逐项核对两侧谱。}
\pitfall{把猜想当成已证的定理。$SU(2)$ 的 $\mathcal N=4$ 理论本身至今没有被证明满足这个对偶，KW06 是把 $S$-对偶当作工作假设来用，只在某些拓扑扇区与极限下验证。另一个错法：以为它只适用于 $SU(2)$——对偶必须同时换群，$SU(N)$ 的对偶是 $PSU(N)$。}
\end{word}

% ============================================================
% 3. S-对偶
% ============================================================
\begin{word}{S-对偶}{S-duality}{s-dual}{§1 p.4；§2.2 p.11}
\mathdef{作用在「群加耦合」这一对上的变换：$S$ 把 $(G,\tau)$ 送到 $({}^LG,-1/n_{\mathfrak g}\tau)$，同时把电、磁荷按 $(m\ n)\mapsto(m\ n)M^{-1}$ 重排；它与 $T:\tau\mapsto\tau+1$ 一起生成 $\Gamma\subset SL(2,\mathbb R)$。等价关系：$(G,\tau)\sim({}^LG,-1/n_{\mathfrak g}\tau)$。不变量是 $\Gamma$ 在 $\tau$ 上的轨道。单连通型时 $n_{\mathfrak g}=1$、$\Gamma=SL(2,\mathbb Z)$。}
\mathex{$G=SU(2)$、$\theta=0$：$S:\tau\mapsto-1/\tau$。$\tau=i$（即 $e^2=4\pi$）是不动点，$-1/i=i$，稳定子为 $\mathbb Z_2$；$\tau\to i\infty$（弱耦合）与 $\tau\to0$（强耦合）被 $S$ 互换。$G=G_2$ 时 $n_{\mathfrak g}=3$，$S:\tau\mapsto-1/(3\tau)$，不动点是 $\tau=i/\sqrt3$。}
\physdef{$4$ 维 $\mathcal N=4$ 理论在 $(G,\tau)$ 与 $({}^LG,-1/n_{\mathfrak g}\tau)$ 下等价的精确量子对称。它交换 Wilson 线与 't Hooft 线、交换电与磁荷、并且（在扭化理论里）作用在典范参数 $\Psi$ 上：$\Psi\mapsto-1/n_{\mathfrak g}\Psi$（§3）。二维约化之后它就是镜像对称（§5）。}
\physex{$G=SU(2)$、$\tau=i$：$\theta=0$、$e^2=4\pi$，理论与自身对偶，$S$ 把 $W$ 玻色子送到磁单极子。$\tau=e^{i\pi/3}=1/2+i\sqrt3/2$（$\theta=\pi$、$e^2=8\pi/\sqrt3$）是 $TS$ 的三阶椭圆点，说明 $\theta$ 与 $e$ 不能分开调。}
\bridge{§2.2 p.11：$S$ 的矩阵为 $\begin{pmatrix}0&1\\-1&0\end{pmatrix}$（单连通型）或 $\begin{pmatrix}0&1/\sqrt{n_{\mathfrak g}}\\-\sqrt{n_{\mathfrak g}}&0\end{pmatrix}$；§3 p.20 给出 $\Psi\mapsto-1/n_{\mathfrak g}\Psi$。§5 的核心论断：$S$-对偶在 $\Sigma\times C$ 上紧化后变成 $\mathcal M_H(G)$ 与 $\mathcal M_H({}^LG)$ 的镜像对称。}
\intuit{S-对偶是强--弱耦合的桥：它让算不动的强耦合理论变成一个弱耦合理论的等价描述。这个类比靠得住，是因为 $\tau$ 住在上半平面、超对称使 $\Gamma$ 成为作用在 $\tau$ 上的离散群，而 $S$ 正是其中那个反演变换。}
\pitfall{把 S-对偶说成「$g\mapsto1/g$」。在 KW06 的归一化下正确形式是 $\tau\mapsto-1/n_{\mathfrak g}\tau$：它同时改 $\theta$（$S$ 不是 $\theta$ 的简单反号）、带 lacing 数、并且换群。只写 $g\mapsto1/g$ 会漏掉 $G\mapsto{}^LG$，从而把 $SU(2)$ 与 $SO(3)$ 的谱搞混。}
\end{word}

% ============================================================
% 4. T-对偶
% ============================================================
\begin{word}{T-对偶}{T-duality}{t-dual}{§1 p.3}
\mathdef{$2$ 维 sigma 模型沿靶空间环面纤维的半径反演：靶空间含环面 $T^d=\mathbb R^d/\Lambda$ 时，把 $\Lambda$ 换成对偶格 $\Lambda^*$，$2$ 维理论的谱不变。数学上它是「对偶阿贝尔簇」的互换：环面 $A$ 与其对偶 $A^\vee=\mathrm{Pic}^0(A)$ 给出同一个 sigma 模型，实现方式是 Fourier--Mukai 变换。不变量是 $2$ 维共形场论的配分函数，动量模与绕数模在谱里对称出现。}
\mathex{取自由玻色子紧化在半径 $R$ 的圆上：$M^2=(m/R)^2+(nR/\alpha')^2$，$m$ 是动量、$n$ 是绕数；$R\mapsto\alpha'/R$ 交换二者，$R=\sqrt{\alpha'}$ 为自对偶点。复几何版本：椭圆曲线 $E=\mathbb C/(\mathbb Z+\tau\mathbb Z)$ 与其对偶 $\mathrm{Pic}^0(E)\cong E$，模参数 $\tau\mapsto-1/\tau$。}
\physdef{$2$ 维超对称 sigma 模型的半径反演对称。它不作用于 $4$ 维规范理论本身；KW06 的用法是：把 $4$ 维理论紧化在 $C$ 上得到靶空间为 $\mathcal M_H$ 的 $2$ 维 sigma 模型，而四维的 $S$-对偶在二维表现为 Hitchin 纤维化纤维上的 T-对偶。所谓「$S$-对偶约化为 T-对偶」即指此。}
\physex{$C$ 取椭圆曲线：$\mathcal M_H(U(1),C)=T^*\mathrm{Jac}(C)$，Hitchin 纤维就是 $\mathrm{Jac}(C)\cong C$；$S$ 诱导的 T-对偶把纤维 $C$ 换成对偶环面 $\mathrm{Pic}^0(C)$。对一般的 $G$，同一个 $b\in B$ 上两条纤维 $F_b$ 与 ${}^LF_b$ 互为对偶环面，$S$ 在它们之间做纤维反演。}
\bridge{§1 p.3：Hitchin 模空间的 $T$-对偶在数学上被理解为 Fourier--Mukai 变换，是几何 Langlands 的某种半经典近似（原文引文献 22、23、24）。§5：Hausel--Thaddeus 证明 ${}^LF_b$ 上的平坦 $U(1)$ 丛与 $F_b$ 的点一一对应，这就是 SYZ 图景在此例中的实现。}
\intuit{直觉是「同一个圆，量半径还是量绕数」：$R$ 与 $1/R$ 描述同一套谱，因为弦既能绕圈也能带动量。这个直觉靠得住，是因为 $2$ 维 CFT 的配分函数在 $R\leftrightarrow1/R$ 下逐项不变，而不是只在低能近似下相等。}
\pitfall{以为四维理论也有 T-对偶。四维没有可反演的半径，T-对偶只出现在紧化到二维之后（或弦的紧化方向上）；KW06 里它是 $S$-对偶在二维的影子，不是四维对偶本身。另一个错法：把 T-对偶当成 $\mathcal M_H$ 的等距——它只保持 $J$ 与 $\omega_K$ 的特定组合，不保持超凯勒度量。}
\end{word}

% ============================================================
% 5. GNO 双群
% ============================================================
\begin{word}{GNO 双群}{GNO dual group}{gno-dual-group}{§1 p.3}
\mathdef{由电荷格的对偶定出的群：$G$ 的电荷取值于权格 $\Lambda_{\mathrm{wt}}(G)=X^*(T)$，Dirac 量子化要求磁荷取值于对偶格 $X_*(T)=\Lambda_{\mathrm{cwt}}(G)$；而半单群由根数据 $(\Lambda_{\mathrm{wt}},\Lambda_{\mathrm{cwt}},\Delta,\Delta^\vee)$ 唯一决定，故「以 $\Lambda_{\mathrm{cwt}}(G)$ 为权格、以 $\Lambda_{\mathrm{wt}}(G)$ 为余权格的群」有确定含义，它就是 GNO 双群。等价关系：双群的双群回到 $G$。}
\mathex{$G=U(1)$：$\Lambda_{\mathrm{wt}}=\Lambda_{\mathrm{cwt}}=\mathbb Z$，双群是 $U(1)$ 自己。$G=SU(2)$：两个格都同构于 $\mathbb Z$，生成元分别是基本权 $\varpi$ 与余根 $\alpha^\vee$，双群是 $SO(3)=SU(2)/\mathbb Z_2$；$G=SO(3)$ 时反过来是 $SU(2)$；$G=SU(3)$ 的双群是 $PSU(3)$。}
\physdef{1976 年 Goddard--Nuyts--Olive 的观察：规范理论中电荷取值于规范群的权格，磁荷取值于某个对偶群的权格，这个对偶群就是 GNO 双群。它是磁荷谱的群论框架：'t Hooft--Polyakov 单极子的磁荷是 $G$ 的一个余权（模 Weyl 群），而余权格正是双群的权格。}
\physex{$G=SU(2)$ 破缺到 $U(1)$：单极子的磁荷为 $\alpha^\vee$（记作 1），$W$ 玻色子的电荷为 $\varpi$（记作 1），Dirac 配对 $\langle\varpi,\alpha^\vee\rangle=1$ 使电荷与磁荷之积量子化为整数。换成 $G=SO(3)$：电荷格缩为 $2\mathbb Z\varpi$、磁荷格扩大一倍，于是多出一条 $w_2=1$ 的 't Hooft 线。}
\bridge{§1 p.3：KW06 明确说「对规范群 $G$，GNO 双群实际上就是 Langlands 对偶群 ${}^LG$」，并归功于 Atiyah 1977 年的提示。出发点是格还是根数据并不重要：磁荷格等于 ${}^LG$ 的权格、也等于 $G$ 的余权格。这条等式是 §6.2 用 ${}^LG$ 的不可约表示给 't Hooft 算子贴标签的依据。}
\intuit{GNO 对偶的直觉是「两个取值格对调，群就跟着对调」：电荷住一个格、磁荷住对偶格，把两格的角色互换就得到另一个理论。这个直觉靠得住，是因为半单群由根数据唯一决定，而对调格恰好就是对偶化根数据的全部内容。}
\pitfall{把 GNO 双群当成「与 $G$ 同构的另一个群」，或以为它只由 Lie 代数决定。它由格决定，因而依赖整体形式：$SU(2)$ 与 $SO(3)$ 有同一个 Lie 代数，双群却互换。另一个错法是把磁荷说成取值于 $G$ 的权格——那是电荷；磁荷取的是余权格，即 ${}^LG$ 的权格。}
\end{word}

% ============================================================
% 6. Langlands 对偶群
% ============================================================
\begin{word}{Langlands 对偶群}{Langlands dual group}{langlands-dual}{§1 p.3}
\mathdef{设 $G$ 是分裂约化群，取极大环面 $T$，根数据为 $(X^*(T),X_*(T),\Delta,\Delta^\vee)$；交换两个格、交换单根与余单根，得 $(X_*(T),X^*(T),\Delta^\vee,\Delta)$，它定义另一个分裂约化群 ${}^LG$。等价关系：${}^{LL}G\cong G$。不变量：Lie 代数一般不同构（$B_n$ 与 $C_n$ 互换；$G_2,F_4,E_8$ 自对偶是例外），而中心与基本群互换，即 $\pi_1(G)\cong Z({}^LG)$、$Z(G)\cong\pi_1({}^LG)$。}
\mathex{${}^LSU(N)=PSU(N)=SU(N)/\mathbb Z_N$；${}^LSp(n)=SO(2n+1)$；${}^LSpin(2n)=SO(2n)/\mathbb Z_2$；${}^LE_8=E_8$、${}^LG_2=G_2$（自对偶）；${}^LU(N)=U(N)$。$N=2$ 时 $SU(2)$ 与 $SO(3)$ 互换。}
\physdef{物理上它是 $S$-对偶交换的那个规范群：$(G,\tau)$ 与 $({}^LG,-1/n_{\mathfrak g}\tau)$ 描述同一个理论。它的不可约表示给 $G$ 的 't Hooft 算子贴标签（§6.2）；它的 Hitchin 模空间 $\mathcal M_H({}^LG,C)$ 是 $\mathcal M_H(G,C)$ 的镜像伙伴（§5）。}
\physex{$G=SU(2)$、$e^2=4\pi$、$\theta=0$：$\tau=i$ 是自对偶点，$SU(2)$ 与 $SO(3)$ 给出等价描述。偏离 $\tau=i$ 后两侧耦合不同，但谱仍互换：$SU(2)$ 的 Wilson 线（按 $SU(2)$ 的表示标号）对偶到 $SO(3)$ 的 't Hooft 线，反之亦然。}
\bridge{§1 p.3--4 与 §2.2 p.11：数学的 ${}^LG$ 与物理的对偶群是同一个构造。全书按这条规则读对偶对：$G$ 换 ${}^LG$、$\Lambda_{\mathrm{wt}}$ 换 $\Lambda_{\mathrm{cwt}}$、$\tau$ 换 $-1/n_{\mathfrak g}\tau$；只有 Lie 代数不同构的 $Sp(k)$ 与 $Spin(2k+1)$ 要另取两个耦合 $\tau$ 与 $\tau'$。}
\intuit{Langlands 对偶群的直觉是「把根与余根对调」：根系记录了长短根，对调就把长短互换，于是群的表示论整体翻到另一侧。这个直觉靠得住，是因为根数据唯一决定分裂约化群，而对调是一次对合，取两次对偶回到原群。}
\pitfall{最常见的错：以为 $SU(N)$ 的对偶是 $SU(N)$ 自己。两者 Lie 代数相同但整体形式不同（$SU(N)$ 单连通、$PSU(N)$ 有 $\mathbb Z_N$ 中心），按根数据的定义对偶群是 $PSU(N)$。另一个错法：以为对偶群的中心与 $G$ 的中心相同——它们恰好互换。}
\end{word}

% ============================================================
% 7. 单连通型（simply-laced）
% ============================================================
\begin{word}{单连通型}{simply-laced}{simply-laced}{§2.2 p.11}
\mathdef{单连通型（simply-laced）：单李代数中所有根等长的情形，即 $A_n$、$D_n$、$E_6,E_7,E_8$ 型；等价地根格与余根格在适当标度下重合，lacing 数 $n_{\mathfrak g}=1$。对这样的 $G$ 有 $S:\tau\mapsto-1/\tau$ 与 $\Gamma=SL(2,\mathbb Z)$。其余情形有长短根：$B_n,C_n$ 与 $F_4$ 的 $n_{\mathfrak g}=2$，$G_2$ 的 $n_{\mathfrak g}=3$。}
\mathex{$G=SU(3)$（$A_2$）：根等长，$n_{\mathfrak g}=1$，$S:\tau\mapsto-1/\tau$，$\Gamma=SL(2,\mathbb Z)$，且 ${}^LSU(3)=PSU(3)$。对照 $G=Sp(2)$（$C_2$）：$n_{\mathfrak g}=2$，$S:\tau\mapsto-1/(2\tau)$，${}^LSp(2)=SO(5)$；$G=G_2$：$n_{\mathfrak g}=3$，$S:\tau\mapsto-1/(3\tau)$。}
\physdef{它决定 $S$-对偶的形式：$n_{\mathfrak g}=1$ 时 $\tau\mapsto-1/\tau$ 与 $\tau\mapsto\tau+1$ 都是对称，$\Gamma=SL(2,\mathbb Z)$；否则 $S$ 带 $\sqrt{n_{\mathfrak g}}$，$\Gamma$ 变成 Hecke 群（§2.2 p.11）。物理上 $n_{\mathfrak g}$ 来自 BPS 质量公式中根与余根的长度比，也能从弦论实现的归一化看出。}
\physex{$G=G_2$（非单连通型，$n_{\mathfrak g}=3$）：$S$ 的不动点是 $\tau=i/\sqrt3$，因为 $-1/(3\cdot i/\sqrt3)=i/\sqrt3$；它有 $|\tau|=1/\sqrt3$，不是模群的椭圆点。$G=SU(2)$（$A_1$，单连通型）：$S$ 的不动点是 $\tau=i$，落在经典的模群图景里。}
\bridge{§2.2 p.11：原文先写单连通型时 $S=\begin{pmatrix}0&1\\-1&0\end{pmatrix}$ 并生成 $SL(2,\mathbb Z)$，再写一般情形的 $S=\begin{pmatrix}0&1/\sqrt{n_{\mathfrak g}}\\-\sqrt{n_{\mathfrak g}}&0\end{pmatrix}$。这条分界线在 §3 决定完整对偶群何时是 $SL(2,\mathbb Z)$、何时是 Hecke 群；§6.2 末又用它判断 $\Psi$ 的有理值能否都由 $\Psi=0,\infty$ 生成。}
\intuit{直觉就是「根都一边长」：没有长短根的差别，根与余根可以等比例对齐，于是 $S$ 只需要 $\tau\mapsto-1/\tau$ 这一步。这个直觉靠得住，是因为 $B_n/C_n$ 的长短根比给出 $n_{\mathfrak g}=2$、$G_2$ 给出 3，恰好就是 $S$ 矩阵里 $\sqrt{n_{\mathfrak g}}$ 的来源。}
\pitfall{把「单连通型」按字面读成「单连通」（$\pi_1(G)=1$），这是两个无关的概念：$PSU(3)$ 是单连通型但 $\pi_1=\mathbb Z_3$；$Sp(2)=Spin(5)$ 单连通却属于 $C_2$、不是单连通型。混淆会直接错判 $n_{\mathfrak g}$，把 $G_2$ 的 $S$ 写成 $\tau\mapsto-1/\tau$（应为 $-1/(3\tau)$）。}
\end{word}

% ============================================================
% 8. 对偶群 Γ
% ============================================================
\begin{word}{对偶群 $\Gamma$}{duality group}{dual-group-gamma}{§2.2 p.11}
\mathdef{由 $S$ 与 $T$ 生成的 $SL(2,\mathbb R)$ 的离散子群 $\Gamma$：$T:\tau\mapsto\tau+1$、$S:\tau\mapsto-1/n_{\mathfrak g}\tau$，按 $\tau\mapsto(a\tau+b)/(c\tau+d)$ 作用在上半平面。它不是规范群，而是作用在耦合 $\tau$（以及 §3 的典范参数 $\Psi$）上的矩阵群。单连通型时 $\Gamma=SL(2,\mathbb Z)$；否则是 $SL(2,\mathbb R)$ 的无限离散子群，即 Hecke 群。不变量是 $\Gamma$-轨道。}
\mathex{$G=SU(2)$：$\Gamma=SL(2,\mathbb Z)$，$\tau=i\infty$（弱耦合）与 $\tau=0$（强耦合）被 $S$ 互换，$\tau=i$ 的稳定子是 $\mathbb Z_2$。$G=G_2$：$S=\begin{pmatrix}0&1/\sqrt3\\-\sqrt3&0\end{pmatrix}$ 的元素不是整数，故 $\Gamma$ 不是 $SL(2,\mathbb Z)$ 的子群。}
\physdef{对偶群是「理论对称性在耦合空间上的像」：每个 $\gamma\in\Gamma$ 提升为作用在 $\mathcal N=4$ 理论上的对称，同时还在超荷代数上诱导一个自同构（§2.2）。$T$ 是经典的（$\theta$ 平移 $2\pi$），$S$ 是量子的（强--弱）。四维流形上的配分函数 $Z_\xi$ 按 $\Gamma$ 的幺正表示变换（§2.2 末、§7）。}
\physex{$G=SU(2)$、$\tau=i$：$S$ 是不动点，$T$ 把 $\theta$ 平移 $2\pi$ 后回到等价理论，$\Gamma$ 在 $i$ 处的轨道给出模曲线上的椭圆点。$G=G_2$：$\tau=i/\sqrt3$ 是 $S$ 的不动点，$\Gamma$ 的基本域比模群的更窄，这正是 §3 里 $\Psi$ 只能落在 $\Gamma$-轨道上的原因。}
\bridge{§2.2 p.11：$\Gamma$ 由 $S,T$ 生成，是 $SL(2,\mathbb R)$ 的无限离散子群；原文强调非单连通型时它不是 $SL(2,\mathbb Z)$ 本身。§3 p.20：$\Gamma$ 作用在 $\Psi$ 上，$\Psi\mapsto(a\Psi+b)/(c\Psi+d)$；§6.2 末：$\Psi$ 的所有有理值恰好是 $\Psi=0,\infty$ 的 $\Gamma$-轨道。}
\intuit{直觉上对偶群是「耦合空间的模群」：它像模群作用在模参数 $\tau$ 上，只是生成元 $S$ 带了 $\sqrt{n_{\mathfrak g}}$。这个直觉靠得住，是因为量子修正只依赖 $\tau$ 的 $\Gamma$-轨道，而 $\Gamma$ 正是由 $T$（经典）与 $S$（量子）实际生成的那个群。}
\pitfall{把对偶群 $\Gamma$ 与规范群 $G$ 混为一谈。$\Gamma$ 是作用在 $\tau$ 上的 $2\times2$ 矩阵群，只在 $n_{\mathfrak g}$ 这个数上依赖 $G$。另一个错法：以为 $\Gamma$ 总是 $SL(2,\mathbb Z)$——非单连通型时它更小，而且 $\Gamma$ 不是 $SL(2,\mathbb Z)$ 的子群。}
\end{word}

% ============================================================
% 9. theta 角
% ============================================================
\begin{word}{$\theta$ 角}{theta angle}{theta-angle}{§1 p.3；§2.2 p.11}
\mathdef{$\theta$ 是拓扑项的系数：作用量含 $\frac{i\theta}{8\pi^2}\int\operatorname{Tr}F\wedge F$。对单连通 $G$，$k=\frac{1}{8\pi^2}\int\operatorname{Tr}F\wedge F\in\mathbb Z$ 是瞬子数（第二 Chern 数），只有 $e^{i\theta k}$ 进入配分函数，故 $\theta$ 是周期 $2\pi$ 的角变量，即 $\theta\in\mathbb R/2\pi\mathbb Z$。数学上它参数化 $H^4(BG,\mathbb Z)\cong\mathbb Z$ 上的特征 $k\mapsto e^{i\theta k}$，配分函数 $Z(\theta)=\sum_kZ_ke^{i\theta k}$ 是 $e^{i\theta}$ 的傅里叶级数。}
\mathex{$G=SU(2)$、$\mathbb R^4$ 上的 $k=1$ BPST 瞬子：$\frac{1}{8\pi^2}\int\operatorname{Tr}F\wedge F=1$，其扇区权重为 $e^{i\theta}e^{-8\pi^2/e^2}$。$\theta=\pi$ 时 $k$ 与 $-k$ 的权重反号，给出符号交替；$\theta=0$ 时两者同号相加。}
\physdef{$\theta$ 角是 $\mathcal N=4$ 理论中 $F\wedge F$ 拓扑项的系数，与规范耦合 $e$ 合成复耦合 $\tau=\theta/2\pi+4\pi i/e^2$。它不改变经典运动方程（因为 $F\wedge F$ 是拓扑项），只改变不同瞬子扇区之间的相对相位；$T:\tau\mapsto\tau+1$ 就是 $\theta\mapsto\theta+2\pi$。}
\physex{$G=SU(2)$、$e^2=4\pi$：$\theta=0$ 给 $\tau=i$（$S$-对偶不动点），$\theta=\pi$ 给 $\tau=1/2+i$。后者在 $S$ 下变成 $-1/(1/2+i)=-2/5+4i/5$，不再停在 $\Re\tau=1/2$ 上——这正是「$\theta$ 不能脱离 $\tau$ 单独谈论」的具体表现。}
\bridge{§1 p.3：$\theta$ 与 $e$ 一起给出 $\tau$；§2.2 p.11：$T:\tau\mapsto\tau+1$ 是 $\theta$ 的 $2\pi$ 平移，对任何 $G$ 都是经典对称。§3 p.20：$\theta$ 又与扭化参数 $t$ 合成典范参数 $\Psi$——所以在几何 Langlands 里真正起作用的旋钮是 $\Psi$，不是 $\theta$ 本身。}
\intuit{$\theta$ 角的直觉是「给瞬子数配一个相位」：每个瞬子扇区贡献 $e^{i\theta k}$，像绕数上的 Bloch 相位。这个直觉靠得住，是因为 $F\wedge F$ 的积分是拓扑不变量，只依赖瞬子数 $k$，完全不影响局域方程。}
\pitfall{把 $\theta$ 当成可以独立调节的耦合。按字面理解会以为「$\theta$ 与 $e$ 是两个独立旋钮」，但 $S$-对偶把两者混在一起（$\tau\mapsto-1/\tau$ 同时改实部与虚部），只有 $\tau$ 是良定义的参数。另一个错法：以为 $\theta$ 的周期总是 $2\pi$——当 $G$ 不单连通时 $T$ 未必是对称，周期可以更大。}
\end{word}

% ============================================================
% 10. 复耦合 τ
% ============================================================
\begin{word}{复耦合 $\tau$}{complexified coupling}{complex-coupling}{§2.2 p.11}
\mathdef{复耦合 $\tau=\theta/2\pi+4\pi i/e^2$ 住在上半平面 $\mathbb H$，因为 $\Im\tau=4\pi/e^2>0$。它是参数空间 $\mathbb H$ 上的一点，对偶群 $\Gamma$ 按 $\tau\mapsto(a\tau+b)/(c\tau+d)$ 分式线性地作用；等价关系是 $\tau\sim\gamma\tau$（$\gamma\in\Gamma$）给出同一个物理理论。不变量是 $\Gamma$-轨道，即模曲线 $X_\Gamma$ 上的一点。}
\mathex{$G=SU(2)$、$\theta=0$：$\tau=4\pi i/e^2$ 落在虚轴上，$e^2=4\pi$ 时 $\tau=i$。$G=G_2$：$\tau=i/\sqrt3$ 是 $S$ 的不动点。$G=U(1)$：$\tau$ 按模群作用，配分函数在 $\Gamma$ 下按幺正表示变换，变换因子与椭圆曲线模形式的情形同型。}
\physdef{$\tau$ 同时编码规范耦合与 $\theta$ 角：$\Im\tau=4\pi/e^2$ 越大耦合越弱，$\Re\tau=\theta/2\pi$。$S$-对偶把它换成 $-1/n_{\mathfrak g}\tau$，$T$ 换成 $\tau+1$。$\mathcal N=4$ 的超对称把量子修正全部吸收进 $\tau$，使它成为理论唯一的连续参数（不计 Higgs 分支的模）。}
\physex{$G=SU(2)$、$e=2$、$\theta=\pi$：$\tau=1/2+\pi i$。$S$ 给出 $-1/\tau=(-1/2+\pi i)/(1/4+\pi^2)$，虚部从 $\pi$ 降到约 $0.31$，即耦合变弱，同时规范群换成 $SO(3)$——一次 $S$ 变换同时完成了强--弱互换与换群。}
\bridge{§2.2 p.11：$\tau=\theta/2\pi+4\pi i/e^2$，$S$ 与 $T$ 生成 $\Gamma$。§3 p.20 的典范参数 $\Psi$ 与 $\tau$ 不同：$\Psi=\Re\tau+i\Im\tau\cdot(t-t^{-1})/(t+t^{-1})$ 还含扭化参数 $t$，$t=\pm i$ 时 $\Psi=\infty$；二维层面真正起作用的是 $\Psi$。}
\intuit{$\tau$ 像复结构的模参数：虚部管耦合强度、实部管 $\theta$ 相位，上半平面正是耦合为正的区域。这个直觉靠得住，是因为作用量里动能项系数正比于 $\Im\tau$、拓扑项系数正比于 $\Re\tau$，两者以 $\tau$ 的形式合成同一个参数。}
\pitfall{两个常见混淆。第一，把 $\tau$ 与曲线 $C$ 的模参数（传统上也记 $\tau$）当成一回事：前者是规范耦合住在上半平面，后者是黎曼面的复结构。第二，把 $\tau$ 与典范参数 $\Psi$ 当成同一个东西：$\Psi$ 还含扭化参数 $t$，并且可以取 $\infty$（当 $t=\pm i$ 时），而 $\tau$ 永远在 $\mathbb H$ 里。}
\end{word}

% ============================================================
% 11. 权格
% ============================================================
\begin{word}{权格}{weight lattice}{weight-lattice}{§2.2 p.13}
\mathdef{权格 $\Lambda_{\mathrm{wt}}(G)=X^*(T)=\operatorname{Hom}(T,U(1))$（$T$ 为极大环面）：由所有不可约表示的最高权生成的自由阿贝尔群，含根格 $\Lambda_{\mathrm{rt}}$ 为有限指数子格；也可写成 $\{\lambda\in\mathfrak t^*:\langle\lambda,\alpha^\vee\rangle\in\mathbb Z,\ \forall\alpha\in\Delta\}$。对偶格是余权格，配对 $\Lambda_{\mathrm{wt}}\times\Lambda_{\mathrm{cwt}}\to\mathbb Z$ 完美。不变量：指数 $[\Lambda_{\mathrm{wt}}:\Lambda_{\mathrm{rt}}]=|Z(G)|$（半单时）。}
\mathex{$G=SU(2)$：$\Lambda_{\mathrm{wt}}=\mathbb Z\varpi$（$\varpi=\alpha/2$ 为基本权），根格为 $2\mathbb Z\varpi$，指数 2 等于 $|Z(SU(2))|$。$G=SO(3)$：权格缩为根格 $2\mathbb Z\varpi$，所以只有整数自旋的表示能做 Wilson 线。$G=U(1)$：权格就是 $\mathbb Z$。}
\physdef{电荷取值于权格：Wilson 线 $W_R$ 的电荷是表示 $R$ 的最高权（模 Weyl 群），因此允许的电荷谱完全由 $\Lambda_{\mathrm{wt}}$ 给出。群的整体形式决定哪些表示合法：$SU(2)$ 允许 $\varpi$ 的奇倍（半整数自旋），$SO(3)$ 只允许偶倍（整数自旋）。}
\physex{$G=SU(2)$：Wilson 线按自旋 $j$ 标号，电荷为 $\pm j$（以 $\varpi$ 为单位，$2j\in\mathbb Z$）；$G=SO(3)$ 时只剩 $j\in\mathbb Z$。$S$-对偶与此对应：$SU(2)$ 的磁荷谱（余权格）恰好等于 $SO(3)$ 的电荷谱，反之亦然。}
\bridge{§2.2 p.13：$\vec n$（电）取值于 $\mathfrak g$ 的权格、$\vec m$（磁）取值于余权格。GNO 的原始观察「磁荷取对偶群的权格」正落在这条等式上：$\Lambda_{\mathrm{wt}}({}^LG)=X_*(T)=\Lambda_{\mathrm{cwt}}(G)$。Langlands 对偶的根数据对调，落地时就是这两个格的互换。}
\intuit{权格是「表示论能给出的全部电荷」：每个表示的最高权都落在这个格上。这个直觉靠得住，是因为权格按定义由表示生成，而电荷就是表示的最高权；根格则是规范场自身（伴随表示）携带的电荷，所以两者的指数恰好是中心的大小。}
\pitfall{把权格与根格混为一谈。对 $SU(2)$ 两者差因子 2：$W$ 玻色子（伴随表示）的电荷在根格 $2\mathbb Z\varpi$ 里，而基本表示的电荷是 $\varpi$，不在根格中。这个差别正是 $SU(2)$ 与 $SO(3)$ 有不同 Wilson 线谱的原因。另一个错法：以为磁荷也在权格里——磁荷在余权格。}
\end{word}

% ============================================================
% 12. 余权格
% ============================================================
\begin{word}{余权格}{coweight lattice}{coweight-lattice}{§2.2 p.13}
\mathdef{余权格 $\Lambda_{\mathrm{cwt}}(G)=X_*(T)=\operatorname{Hom}(U(1),T)$：由基本余权生成的格，含余根格为有限指数子格（指数为 $|\pi_1(G)|$），也是权格的对偶格：完美配对 $\Lambda_{\mathrm{wt}}\times\Lambda_{\mathrm{cwt}}\to\mathbb Z$。关键恒等式：${}^LG$ 的权格等于 $G$ 的余权格。}
\mathex{$G=SU(2)$：$\Lambda_{\mathrm{cwt}}=\mathbb Z\alpha^\vee$（$\alpha^\vee$ 为余根），与权格同构但住在对偶空间里。$G=SO(3)=SU(2)/\mathbb Z_2$：余权格扩大为 $\frac12\mathbb Z\alpha^\vee$，比 $SU(2)$ 的大一倍——这正是 $SO(3)$ 多出 $w_2=1$ 的 't Hooft 线的来源。}
\physdef{磁荷取值于余权格。't Hooft 线绕其一周的奇异性由一个同态 $U(1)\to G$ 的共轭类决定，也就是一个余权（模 Weyl 群）。换成对偶群的语言：$G$ 的 't Hooft 线由 ${}^LG$ 的权标号，而 ${}^LG$ 的权格正是 $G$ 的余权格。}
\physex{$G=SU(2)$：'t Hooft--Polyakov 单极子的磁荷为 $\alpha^\vee$（记作 1），绕它的 $S^2$ 上 $\int_{S^2}F$ 非零。$G=SO(3)$：同一个 Lie 代数，但磁荷格含 $\frac12\alpha^\vee$，于是存在 $w_2=1$ 的 't Hooft 线，它在 $SU(2)$ 理论里不存在。}
\bridge{§2.2 p.13：$\vec m$ 取值于余权格；§6.2 p.84：'t Hooft 算子 $T({}^LR)$ 由 ${}^LG$ 的最高权标号，而 $\Lambda_{\mathrm{wt}}({}^LG)=\Lambda_{\mathrm{cwt}}(G)$；同处还指出 $S^2$ 上的 $G$-丛由 $\pi_1(G)\cong Z({}^LG)$ 分类，丛非平凡当且仅当中心在表示上作用非平凡。}
\intuit{余权格就是「磁荷格」：单极子的磁荷是把 $U(1)$ 映进 $G$ 的方式，而同态 $U(1)\to T$ 按定义就是一个余权。这个直觉靠得住，是因为绕单极子的 $S^2$ 上的 $G$-丛恰好由这样的同态（模 Weyl 群）分类，即 Grothendieck 对 $\mathbb P^1$ 上 $G$-丛的分类。}
\pitfall{以为余权格与权格「同构所以是同一个格」。$SU(2)$ 与 $SO(3)$ 的 Lie 代数相同，余权格却差一倍，这个差别决定两者 't Hooft 线谱不同。另一个错法：忘了模 Weyl 群——磁荷只在 Weyl 群轨道意义下良定义，把某个具体余权当作「唯一的磁荷」会数错线的条数。}
\end{word}

% ============================================================
% 13. GNO 线
% ============================================================
\begin{word}{GNO 线}{GNO line}{gno-line}{§6.2 p.84}
\mathdef{给定 $G$ 的一个不可约表示 ${}^LR$（等价地，${}^LG$ 的一个最高权 $\lambda$）与一个一维流形 $S$，'t Hooft 线算子记作 $T({}^LR,S)$。它的定义是：沿 $S$ 制造一个余维 2 的奇异，使绕 $S$ 的 $S^2$ 上 $G$-丛的拓扑由同态 $\rho:U(1)\to G$ 的共轭类给出，而该共轭类恰好由 $\lambda$ 分类。因此「线的标号」是 ${}^LG$ 的一个表示，而不是 $G$ 的表示。}
\mathex{奇异性形式（原文 (6.9)）：$F=\frac{i}{2}\star_3 d\frac{1}{|\vec x|}$、$\phi=\frac{i}{2|\vec x|}ds$，即绕线的 $S^2$ 上 $\int F$ 非零。$\rho$ 的共轭类由 ${}^LG$ 的最高权分类，因为 $\operatorname{Hom}(U(1),G)$ 的共轭类与 ${}^LG$ 的权一一对应（两者都模 Weyl 群）。}
\physdef{在四维 $\mathcal N=4$ 规范理论里，Wilson 线按 $G$ 的表示标号、't Hooft 线按 ${}^LG$ 的表示标号，两者住在同一套对偶框架里。当 $G$ 有非平凡中心时，'t Hooft 线还带一个拓扑类（$S^2$ 上 $G$-丛的 $w_2$）：$T({}^LR)$ 的丛非平凡当且仅当 $Z({}^LG)$ 在 ${}^LR$ 上作用非平凡。}
\physex{$G=SO(3)$：'t Hooft 线按 ${}^LSO(3)=SU(2)$ 的表示标号，即自旋 $j$；当 $j$ 为半整数时 $Z_2=Z(SU(2))$ 作用非平凡，对应 $SO(3)$ 里 $w_2=1$ 的 't Hooft 线，这种线在 $SU(2)$ 理论里没有（$\pi_1(SU(2))=1$）。$G=PSU(N)$：'t Hooft 线按 ${}^LPSU(N)=SU(N)$ 的表示标号，即按 $N$ 元性取标签。}
\bridge{§6.2 p.84：「$G$ 的 't Hooft 算子由 ${}^LG$ 的表示分类」就是经典 GNO 对偶，近来被重解为「关于算子而非态的对偶」。GNO 线的标号规律与 §2.2 p.13 的磁荷格是同一件事：磁荷取 $G$ 的余权格 $=$ ${}^LG$ 的权格 $=$ ${}^LG$ 表示的标签。}
\intuit{GNO 线的直觉是「沿一条线挖掉一块看丛的绕数」：绕线的 $S^2$ 上丛由中心的数据刻画，而这份数据正好是 ${}^LG$ 的一个表示。这个直觉靠得住，是因为 $S^2$ 上的 $G$-丛由 $\pi_1(G)$ 分类、$\pi_1(G)\cong Z({}^LG)$，中心作用在哪个表示上就决定了线的拓扑类。}
\pitfall{把 GNO 线当成几何意义上的「线」，以为它像 Wilson 线一样就是沿 $S$ 的一条世界线。按字面理解会漏掉关键：它其实是「一个标签」（${}^LG$ 的表示）加上一个流形 $S$，拓扑类由中心作用决定；同样一条几何曲线可以有多个不同标号的算子。另一个错法：以为 't Hooft 线按 $G$ 自己的表示标号——那是 Wilson 线。}
\end{word}

% ============================================================
% 14. Hecke 群
% ============================================================
\begin{word}{Hecke 群}{Hecke group}{hecke-group}{§2.2 p.11}
\mathdef{对偶群 $\Gamma$ 由 $T:\tau\mapsto\tau+1$ 与 $S:\tau\mapsto-1/(n_{\mathfrak g}\tau)$ 生成。令 $z=\sqrt{n_{\mathfrak g}}\,\tau$，则 $T$ 变为 $z\mapsto z+\sqrt{n_{\mathfrak g}}$、$S$ 变为 $z\mapsto-1/z$；于是 $\Gamma$ 共轭于 $SL(2,\mathbb R)$ 的离散子群 $\langle z\mapsto z+\sqrt{n_{\mathfrak g}},\ z\mapsto-1/z\rangle$。$n_{\mathfrak g}=1$ 时它是 $SL(2,\mathbb Z)$；$n_{\mathfrak g}=2$ 时是经典 Hecke 群 $H(\sqrt2)$，$n_{\mathfrak g}=3$ 时是 $H(\sqrt3)$——这两者适当共轭后都是模群的有限指数子群，故仍是算术群。}
\mathex{$n_{\mathfrak g}=2$ 的情形：$G=Sp(k)$、$Spin(2k+1)$ 或 $F_4$；$n_{\mathfrak g}=3$：$G=G_2$。$G_2$ 时 $S$ 的矩阵为 $\begin{pmatrix}0&1/\sqrt3\\-\sqrt3&0\end{pmatrix}$，作用为 $\tau\mapsto-1/(3\tau)$，不动点 $\tau=i/\sqrt3$。$SU(2)$ 时 $\Gamma=SL(2,\mathbb Z)$，基本域即模群的基本域。}
\physdef{$\Gamma$ 是「完整对偶群」：它同时含 $T$（$\theta$ 平移 $2\pi$）与 $S$（强--弱对偶），只有 $G$ 单连通型（根全等长）时 $\Gamma$ 才等于 $SL(2,\mathbb Z)$。非单连通型时 $S$ 的矩阵含因子 $1/\sqrt{n_{\mathfrak g}}$、不是整数矩阵，所以 $\Gamma$ 既不是 $SL(2,\mathbb Z)$ 的子群也不是它的超群：两者只在由 $T$ 生成的子群上重合。}
\physex{$G=G_2$：$n_{\mathfrak g}=3$，$\Gamma$ 由 $\tau\mapsto\tau+1$ 与 $\tau\mapsto-1/(3\tau)$ 生成；$S$ 的不动点 $\tau=i/\sqrt3$ 满足 $-1/(3\tau)=\tau$。原文 §6.2 末指出：$\Psi$ 的全部有理值（也只有有理值）都能从 $\Psi=0$ 或 $\Psi=\infty$ 经 $\Gamma$ 变换得到，无论 $\Gamma=SL(2,\mathbb Z)$ 还是 Hecke 群。}
\bridge{§2.2 p.11：非单连通型时 $S:\tau\mapsto-1/(n_{\mathfrak g}\tau)$，生成的是 Hecke 群而不是 $SL(2,\mathbb Z)$；$n_{\mathfrak g}$ 反映根与余根的长度比，$F_4$ 取 2、$G_2$ 取 3，其余取 1。§3 p.20：$\Gamma$ 同样作用在 $\Psi$ 上（$\Psi\mapsto(a\Psi+b)/(c\Psi+d)$），基本对偶是 $\Psi\mapsto-1/(n_{\mathfrak g}\Psi)$。}
\intuit{Hecke 群是「带刻度的模群」：$T$ 的步长仍是平移，但 $S$ 的刻度被 $\sqrt{n_{\mathfrak g}}$ 拉伸，于是基本域的宽度不再是 1 而是 $\sqrt{n_{\mathfrak g}}$。这个直觉靠得住，是因为把 $\tau$ 换元成 $z=\sqrt{n_{\mathfrak g}}\,\tau$ 之后，$S$ 就变回标准的 $z\mapsto-1/z$，整个群由平移 $\sqrt{n_{\mathfrak g}}$ 与反演生成。}
\pitfall{把 Hecke 群与 Hecke 算子（Hecke 特征层）混为一谈：前者是作用在耦合 $\tau$ 或 $\Psi$ 上的对偶群，后者是作用在 D-模上的算子，两者只是共用「Hecke」这个名字。另一个错法：以为所有 $G$ 的对偶群都是 $SL(2,\mathbb Z)$——只有单连通型如此，非单连通型的 $n_{\mathfrak g}=2,3$ 使 $S$ 不再是整数矩阵。}
\end{word}
